\documentclass[11pt]{article}

\usepackage[margin=1in]{geometry}
\usepackage[T1]{fontenc}
\usepackage[utf8]{inputenc}
\usepackage{microtype}
\usepackage{hyperref}
\usepackage{enumitem}
\usepackage{booktabs}
\usepackage{amsmath}

\setlength{\emergencystretch}{2em}

\hypersetup{
  colorlinks=true,
  linkcolor=black,
  citecolor=black,
  urlcolor=black
}

\title{Facts, Trust, and Cooperation:\\
The KFD Foundation Triad}
\author{Keren Dong\\Kungfu Origin Technology Limited\\\href{mailto:keren.dong@kungfu.link}{keren.dong@kungfu.link}}
\date{Draft: July 2026}

\begin{document}

\maketitle

\begin{abstract}
Persistent work must remain intelligible across participants, tools, and time.
That requirement is often described with desirable values such as truth,
trust, and cooperation, but a list of values does not explain what makes a
shared work world possible. This paper presents the first three Kung Fu
Decisions (KFDs) as an ordered systems model rather than a code of conduct:
facts must not drift; trust must start from facts; and cooperation must start
from trusted value. The three decisions preserve, respectively, continuity of
the work world, warranted reliance on that world, and the capacity of
independent participants to generate further joint action. Their composition
is recursive: cooperation produces new occurrences and consequences, which
must become inspectable inputs to the next fact cut. We test the model through
removal and order-inversion arguments and through one cross-session
continuation case. We then show how declared perspective, governed primitive
discovery, and causal experience arise as the first pressures on this
foundation, providing a narrow bridge to KFD-4, KFD-5, and the draft KFD-6.
The contribution is a falsifiable account of the minimal roles required for a
durable cooperative world, not a claim of complete ontology, moral consensus,
or universal sufficiency.
\end{abstract}

\section{Introduction}

A conversation can end while the work continues. A maintainer leaves, an
agent process stops, a repository moves, a release is promoted, or a new
participant arrives with no access to the private context that made the last
decision seem obvious. If the work is to survive, the newcomer must recover
more than a summary. The newcomer must be able to identify what happened,
decide which claims may be relied upon, and determine whether and how to
continue.

This is a systems problem before it is an interface problem. Better models,
larger context windows, and more persuasive explanations do not by themselves
preserve the identity of the work across time. Nor does a complete event log
automatically tell a participant what a claim proves or why cooperation is
worth entering. Persistent work needs a shared world that can outlive the
participant who currently interprets it.

The first three Kung Fu Decisions (KFDs) state a compact foundation for that
world \cite{kfd2026alpha42}:

\begin{quote}
KFD-1: facts must not drift.\\
KFD-2: trust must start from facts.\\
KFD-3: cooperation must start from trusted value.
\end{quote}

Read as three imperatives, the statements can sound like good advice. This
paper advances a stronger and more falsifiable interpretation: they name three
different roles that persistent cooperation cannot collapse into one another.
KFD-1 preserves the continuity of the world about which participants reason.
KFD-2 converts parts of that world into bounded permission to rely. KFD-3
converts warranted reliance and visible value into coordinated action among
participants who retain independent judgment.

The dependency is ordered, but its operation is recursive:

\begin{equation}
\text{facts} \rightarrow \text{trust} \rightarrow \text{cooperation}
\rightarrow \text{new occurrences} \rightarrow \text{successor facts}.
\label{eq:intro-loop}
\end{equation}

This loop is the paper's central object. It explains why the triad can be
foundational without being a complete ontology. The three KFDs do not name
every object a cooperative system will need. They preserve the conditions
under which new action can occur, its consequences can correct the shared
world, and missing objects can later be discovered without replacing reality
with narrative.

The paper makes four contributions:

\begin{itemize}[leftmargin=*]
  \item it defines facts, trust, and cooperation as operational roles rather
        than moral qualities;
  \item it uses removal and order-inversion tests to explain why the roles are
        separately necessary and directionally dependent;
  \item it models the triad as a recursive kernel through which a shared work
        world persists, becomes actionable, and produces further evidence;
        and
  \item it derives a narrow frontier from that kernel: declared perspective
        in KFD-4, separated genesis and qualification in KFD-5, and the
        causal-experience gate in the draft KFD-6.
\end{itemize}

The claims remain bounded. The paper does not prove that the triad is a
mathematically unique basis for every social or technical system. It does not
equate recorded claims with objective reality, trust with certainty, or
cooperation with agreement. It argues that any system claiming durable
cooperation must implement equivalent roles or explain how it preserves the
same continuity, warranted reliance, and generative coordination without
them.

Section~\ref{sec:problem} defines the continuity problem.
Section~\ref{sec:foundation} specifies the three roles.
Section~\ref{sec:necessity} applies removal and order tests.
Section~\ref{sec:kernel} develops the recursive model and a cross-session
case. Section~\ref{sec:frontier} shows how KFD-4/5/6 arise at the first
pressure frontier. Sections~\ref{sec:related} and~\ref{sec:discussion}
position and bound the claim.

\section{Problem: Implicit Boundaries Under Agent Pressure}
\label{sec:problem}

\subsection{Real-world agent work}

We use \emph{real-world agent work} for persistent work in which agent output
can change or support decisions about concrete artifacts, responsibilities, or
external state. Examples include repository changes, builds, package releases,
remote operations, long-running tasks, runtime traces, user approvals, and
recovery after interruption. The work spans several kinds of state:

\begin{itemize}[leftmargin=*]
  \item material state such as source trees, binaries, and documents;
  \item runtime state such as actions, observations, receipts, and failures;
  \item epistemic state such as claims, evidence, gaps, and current decisions;
  \item responsibility state such as purpose, authority, ownership, and next
        action; and
  \item presentation state exposed through human and agent interfaces.
\end{itemize}

No participant sees this entire world directly. Work therefore depends on
representations and handoffs. Distributed cognition already treats cognition
as spread across people and artifacts rather than contained in one head
\cite{hutchins1995cognition}; agent participation increases both the capacity
and the coordination burden of that system.

\subsection{The implicit-boundary pattern}

Let a contact surface be a relation between at least two sides $L=\{l_1,
l_2,\ldots\}$ whose outputs, interpretations, or authorities must meet. Before
the surface becomes load-bearing, coordination may be provided by an implicit
mechanism $H$: memory, convention, reputation, a private prompt, a maintainer's
interpretation, or manual reconstruction.

Pressure changes when a new participant, larger scale, higher frequency,
stronger authority, greater heterogeneity, lower latency tolerance, or larger
consequence enters the system. We denote that change by $\Delta P$. Recurrent
failures $\Phi$ then become visible: translation mistakes, incomplete
handoffs, lost attribution, ambiguous authority, repeated reconstruction, or
unowned responsibility.

The boundary problem can be summarized as:

\begin{equation}
(L, H, \Delta P) \longrightarrow \Phi .
\end{equation}

A primitive candidate $X$ proposes to mediate that surface. It has value only
if it turns part of the old implicit mechanism into a stable, inspectable
object without hiding the responsibility it replaces:

\begin{equation}
X = (I, B, A, Y, O),
\end{equation}

where $I$ is identity, $B$ is boundary, $A$ is authority, $Y$ is lifecycle,
and $O$ is the minimum useful operation set.

\subsection{Three failure classes}

\paragraph{Fact drift.}
The same load-bearing fact appears differently in source configuration,
packaged artifacts, documentation, UI defaults, CLI output, schemas, and
runtime behavior. An agent may read one surface, a human another, and the
runtime enforce a third. More intelligence then accelerates disagreement
instead of resolving it.

\paragraph{Trust by reconstruction.}
A release, task completion, runtime state, or plugin is called trustworthy,
but the supporting facts, evidence cut, residual risk, and responsibility
boundary are unavailable. The recipient must trust authority or repeat the
maintainer's reconstruction. Provenance systems reduce this cost for software
artifacts \cite{slsa,intoto2019}; real-world agent work creates the same need
across action and knowledge boundaries.

\paragraph{Cooperation by capture.}
A product may force participants into a managed workflow, conceal constraints,
or make alternative paths impractical. This can produce compliance while
destroying the participant's ability to inspect value, challenge a claim, or
recover independently. For agents, a private prompt can become the only
interface; for humans, a GUI can become the only place where state appears.

\subsection{Why friction alone is insufficient}

Not every recurring problem deserves a new primitive. The failure may be a bug,
an absent policy, a local abstraction, poor documentation, or a narrow internal
model. Naming a false primitive exports complexity to every consumer. A valid
discovery method must therefore preserve the no-new-primitive alternative and
ask whether deleting the candidate recreates the same coordination burden.
Boundary pressure generates candidates; it does not certify them.

\section{The Foundation Triad}
\label{sec:foundation}

\subsection{KFD-1: continuity of the fact world}

KFD-1 states that facts must not drift. A \emph{fact} here is not whatever a
system happens to store, nor a claim of observer-free certainty. It is a
load-bearing reference admitted from a declared source, authority, and
evidence boundary. KFD distinguishes a \emph{fact cut}, which states what
holds at such a boundary, from a \emph{causal record}, which states what
occurred between declared boundaries \cite{kfd2026alpha42}. Either may include
known gaps, uncertainty, or conflicting evidence.

Let $F_t$ be a fact cut used at time $t$, and let $c$ be its declared
compatibility boundary. If the system presents the same fact identity and
boundary at $t'$ without declaring a change, its load-bearing meaning must be
preserved:

\begin{equation}
\operatorname{id}(F_t)=\operatorname{id}(F_{t'}) \land c_t=c_{t'}
\Longrightarrow
\operatorname{meaning}(F_t)=\operatorname{meaning}(F_{t'}).
\label{eq:fact-continuity}
\end{equation}

The invariant does not freeze the world. It requires semantic change to cross
an explicit boundary so later evidence can contradict, supersede, or narrow an
earlier model without silently rewriting what the earlier participant saw.

KFD-1 therefore provides \emph{world continuity}. It lets two participants
refer to the same admitted state or occurrence even when they disagree about
its interpretation. Without such a reference, disagreement cannot reliably
be distinguished from drift in the object being discussed.

\subsection{KFD-2: warranted reliance}

KFD-2 states that trust must start from facts. Trust is not modeled as a
global property of a person, model, product, or artifact. It is a bounded
permission to rely on a claim for a purpose.

For claim $C$ and purpose $P$, a trust assessment can be represented as:

\begin{equation}
T_P(C)=\mathcal{A}(C,P,F,E,R,Q),
\label{eq:trust-assessment}
\end{equation}

where $F$ is the bound fact cut, $E$ the checked evidence, $R$ residual risk,
and $Q$ the assignment of source, verification, and decision responsibility.
The result may be pass, warning, fail, or unverifiable. Its force is limited to
the stated purpose and cut.

This definition explains why more data is not the same as more trust. A log
may be authentic while irrelevant to a completion claim. A check may pass
while omitting the platform a consumer uses. A capable agent may produce a
persuasive explanation without exposing the facts that would let another
participant challenge it. KFD-2 makes reliance revisable because contrary
facts can change the assessment without first defeating the authority or
reputation of the claimant.

KFD-2 therefore provides \emph{judgment continuity}. It allows a later
participant to recover not only what was claimed, but why and within what
boundary the claim was permitted to guide action.

\subsection{KFD-3: generative coordination}

KFD-3 states that cooperation must start from trusted value. Facts and trust
do not themselves create joint action. A system may expose accurate evidence
while hiding the value exchange, available alternatives, or constraints under
which participation is requested.

Let $V$ be a participant-relevant value claim, $T(V)$ its trust assessment,
$K$ the visible constraints, and $O$ the meaningful available options. A
cooperation decision has the form:

\begin{equation}
C = \mathcal{C}(V,T(V),K,O,Q_C),
\label{eq:cooperation}
\end{equation}

where $Q_C$ declares responsibility for the commitment and resulting action.
The participant need not agree with every other participant. It must be able
to understand what is offered, what is constrained, and what accepting,
rejecting, escalating, or withdrawing means within the declared boundary.

This is not an absence-of-control rule. Sandboxes, permissions, refusal,
revocation, and hard safety gates may be necessary. KFD-3 distinguishes
explainable, reviewable constraint from hidden capture, and cooperation from
mere compliance. The distinction matters operationally: a compliant component
executes inside one authority; a cooperating participant can contribute new
facts, objections, alternatives, and commitments to a shared world.

KFD-3 therefore provides \emph{action continuity}. It lets independent
participants extend the work without surrendering the fact and judgment
boundaries that made the extension trustworthy.

\subsection{Three irreducible questions}

The triad can be summarized by three questions:

\begin{table}[ht]
\centering
\small
\begin{tabular}{p{0.15\linewidth}p{0.24\linewidth}p{0.48\linewidth}}
\toprule
Layer & Role & Question preserved across participants and time \\
\midrule
KFD-1 & World continuity & What state or occurrence are we talking about, and
where may its meaning change? \\
KFD-2 & Judgment continuity & Why may this claim be relied upon for this
purpose, and what remains unproved? \\
KFD-3 & Action continuity & Why and within what visible choices and constraints
should independent participants act together? \\
\bottomrule
\end{tabular}
\caption{The three operational roles of the Foundation Triad.}
\label{tab:triad-roles}
\end{table}

The roles correspond to reality as admitted, belief as warranted, and action
as coordinated. They are related but not interchangeable. That separation is
what makes the foundation testable.

\section{Necessity and Order}
\label{sec:necessity}

The triad is a foundation only if each role carries independent load and if
later roles cannot safely manufacture their own earlier conditions. We apply
two conceptual tests: removal and order inversion.

\subsection{Removal tests}

\paragraph{Remove KFD-1.}
Suppose a system retains trust assessments and cooperation records but does
not preserve stable fact identities and change boundaries. A later participant
cannot determine whether an assessment still refers to the same contract,
artifact, occurrence, or meaning. Trust attaches to a moving object. Conflict
is resolved by the latest narrative or strongest authority because no stable
reference remains against which either can be challenged.

\paragraph{Remove KFD-2.}
Suppose facts are stable and participants are willing to cooperate, but no
purpose-bound assessment connects facts to reliance. Every participant must
either repeat the complete verification, accept another participant's
reputation, or act without knowing which gap matters. The archive may be
excellent, but reliance does not scale. Stable facts are necessary evidence;
they are not self-interpreting permission to act.

\paragraph{Remove KFD-3.}
Suppose facts and assessments are inspectable, but the system has no account of
participant value, choice, or visible constraint. It can support audit,
certification, and command execution. It cannot distinguish durable
cooperation from compliance obtained by monopoly, hidden workflow capture, or
unreviewable pressure. The world may be correct and trustworthy yet remain
non-generative across independent participants.

The failure modes are summarized in Table~\ref{tab:removal}.

\begin{table}[ht]
\centering
\small
\begin{tabular}{p{0.18\linewidth}p{0.31\linewidth}p{0.38\linewidth}}
\toprule
Removed role & What remains & Structural failure \\
\midrule
Stable facts & trust language and coordination & reliance binds to mutable
narrative or authority \\
Warranted trust & records and willing action & repeated verification, blind
reliance, or reputation substitution \\
Trusted-value cooperation & facts and assessment & accountable control may
remain, but independent joint generation collapses into compliance \\
\bottomrule
\end{tabular}
\caption{Removal tests for the three foundation roles.}
\label{tab:removal}
\end{table}

These tests establish functional non-redundancy, not mathematical uniqueness.
Another vocabulary may implement equivalent roles. The KFD claim is that the
capabilities themselves cannot be omitted from durable cooperation.

\subsection{Order-inversion failures}

The arrows in the triad describe dependency, not a rigid wall-clock sequence.
Participants may gather evidence, negotiate value, and act in overlapping
cycles. Nevertheless, a later layer may not retroactively create the basis on
which it claims to stand.

\paragraph{Trust before facts.}
When a system asks for reliance before exposing the relevant facts, identity,
reputation, confidence, or persuasive narrative becomes the substitute. Facts
may later be selected to defend the trust already granted. Correction then
requires defeating an authority relationship rather than updating an
assessment.

\paragraph{Cooperation before trustable value.}
When a system asks for commitment before value, facts, choices, and
constraints can be inspected, participation becomes the evidence that the
offer was acceptable. Lock-in is misread as preference and compliance as
agreement. The system can grow while losing the ability to learn why
participants would have chosen differently.

\paragraph{Successor facts admitted by the action that needs them.}
When completed action is allowed to rewrite its own prior fact cut, a result
can erase failed attempts, costs, reversals, or disconfirming evidence. The
loop becomes self-certifying: success defines which history counts. KFD-1
therefore applies again when new occurrences are admitted into the successor
world.

\subsection{Necessary, not sufficient}

The triad does not guarantee good outcomes. Facts may be incomplete, an
assessment mistaken, visible choices economically unequal, or cooperation
harmful to others outside the declared boundary. Domain knowledge, ethics,
law, security, governance, and empirical correction remain necessary. The
foundation claim is narrower: without continuity of fact, warranted reliance,
and participant-aware coordination, those additional mechanisms cannot form a
durable shared work world.

\section{Three Boundary-Primitive Cases}
\label{sec:evidence}

The three cases in this section arose in different subsystems and have
different maturity. Their common feature is not that they record data. Each
names a contact surface where participants previously relied on implicit human
coordination.

\begin{table}[ht]
\centering
\small
\begin{tabular}{p{0.16\linewidth}p{0.30\linewidth}p{0.38\linewidth}}
\toprule
Candidate & Contact surface & Implicit mechanism replaced \\
\midrule
Xinfa Atlas & project reality $\leftrightarrow$ human/agent understanding &
Repeated repository interpretation, private prompts, hand-written summaries,
and divergent human/agent views \\
Kungfu Episode & agent action $\leftrightarrow$ real-world result &
Sessions, logs, process exit, and human reconstruction of which facts,
artifacts, receipts, and proof belong to one unit of work \\
Release Passport & reviewed source/build $\leftrightarrow$ released artifact &
Maintainer memory about what was promoted, which checks ran, which artifacts
exist, and what users can verify \\
\bottomrule
\end{tabular}
\caption{Three candidate primitives and the implicit boundaries they mediate.}
\label{tab:cases}
\end{table}

\subsection{Xinfa Atlas: reality to shared understanding}

A Xinfa Atlas is defined as the smallest portable, reproducible, and
verifiable semantic closure compiled at a declared source and policy cut, from
which equivalent human and agent views can be derived without reinterpreting
the unbounded source corpus \cite{xinfaAtlas2026}. It binds scope, provenance,
verification state, gaps, conflicts, routes, and deterministic identity while
remaining a representation rather than source authority.

The boundary pressure is cognitive. A repository may be readable to a
maintainer because that person carries years of implicit context. New humans,
agents, documentation sites, IDEs, and task runners otherwise reconstruct
different project models. A context window or archive is too narrow as the
primitive: it names a transport or one reader's selection. Atlas names the
cut-bound closure from which bounded Human Views, Agent Task Charts, and GUI
views can be derived under the same root \cite{xinfaProjections2026}.

The deletion test is strong but not conclusive. Without Atlas, the system
either repeats the full closure definition, elevates a storage pack into the
user object, or allows each projection to become a separate fact source. The
narrower alternative is retrieval over raw repositories. That alternative
would falsify Atlas if it delivered equivalent provenance, portability,
bounded cost, gap visibility, deterministic identity, and cross-reader parity
without recreating Atlas semantics.

As of the evidence cut used by this paper, Atlas identity and bounded
projections are accepted and implemented in Kungfu's Xinfa subsystem. Adoption
and measurable context-cost reduction remain open evaluation questions.

\subsection{Kungfu Episode: action to consequence}

A Kungfu Episode is a bounded causal unit of actual work whose facts,
artifacts, dependencies, receipts, and verification roots can be inspected,
sealed, exported, recovered, and used to support decisions
\cite{kungfuEpisode2026}. The object is deliberately not a chat session,
process, or bag of logs. Those coordinates end for technical reasons, while
real responsibility may span runs and preserve useful partial work after a
failure.

The boundary pressure is causal and accountable. A successful tool call does
not prove visibility, durability, projection, or completion. A process exit
does not state which artifacts belong to the work. A trace can show calls
without defining the authoritative unit that export, fsck, recovery, and later
trust assessment should address. Episode collects those relations under one
bounded lifecycle while keeping Claim distinct from Proof and Decision.

The deletion test asks whether runtime consumers can still export, verify,
recover, and reason about one unit of actual work without reconstructing it
from sessions, logs, traces, and external files. A narrower event-group label
would falsify the need for Episode if it supported the same causal closure,
identity, lifecycle, receipts, and recovery semantics.

Episode is the central product vocabulary and an accepted architectural
object, but its implementation evidence is mixed. Core causal-segment work is
active while some sealed identity, cross-provider equivalence, and bundle
contracts remain proposed or partial. This paper therefore treats Episode as
a strong, actively dogfooded candidate rather than a completed universal
standard.

\subsection{Buildchain Release Passport: build activity to released truth}

Buildchain Release Passport is a durable release responsibility record for
artifacts that users or agents depend on \cite{buildchainPassport2026}. Its
first-read artifact, \texttt{buildchain.release.json}, binds reviewed source,
release line and channel, exact tags, artifact hashes, checks, surface impact,
package publication, residual evidence, and transaction state.

The boundary pressure comes after CI says green. Build systems can compile,
test, and upload while leaving the recipient unable to answer which reviewed
state was released, what exact artifact was produced, which floating channel
now points to it, or how to verify and recover the transaction. Historically a
maintainer reconstructs this story from workflow pages, tags, package
registries, and memory.

The deletion test is whether an independent human or agent can reconstruct the
same release responsibility record from ordinary provider state at comparable
cost and without private maintainer context. The narrower alternative is a
provenance attestation alone. That alternative is sufficient only if it also
captures promotion, compatibility impact, package-set and channel state,
residual risk, and recovery responsibility rather than artifact origin alone.

This is the most mature case in the paper. Buildchain stable release
\texttt{v2.12.7} publishes the passport and its supporting evidence as GitHub
Release assets. Maturity in one ecosystem does not prove generality, but it
does show that the object can bear real release operations rather than remain
only a diagram.

\subsection{Shared structure and bounded claim}

The cases support a common pattern:

\begin{equation}
\text{implicit coordination}
\xrightarrow{\Delta P}
\text{recurring boundary failure}
\xrightarrow{X}
\text{inspectable mediation}.
\end{equation}

They do not yet establish that the pattern is necessary for all primitives or
that any case will reach historical significance. The evidence supports a
candidate-generation method and three non-trivial implementation cases. Time,
independent adoption, and adversarial deletion tests must carry the stronger
claim.

\section{Composition and Evaluation}
\label{sec:evaluation}

The cases become more informative when composed. Composition also reveals why
no one primitive should absorb every authority.

\subsection{A project cut across four dimensions}

Kungfu's Project Cut architecture distinguishes the following objects
\cite{projectCut2026}:

\begin{table}[ht]
\centering
\small
\begin{tabular}{p{0.22\linewidth}p{0.25\linewidth}p{0.39\linewidth}}
\toprule
Object & Dimension & Authority \\
\midrule
Git tree/commit & material/publication & Exact bytes, ancestry, and atomic
publication coordinate \\
Xinfa Atlas & spatial/epistemic & Declared knowledge, evidence state, routes,
gaps, and semantic closure at one cut \\
Kungfu Episode & temporal/causal & Actual actions, facts, artifacts, receipts,
proof inputs, and sealed transition evidence \\
Project Cut & cross-object binding & Versioned references among source,
Atlas, Episode delta, policies, omissions, and roots \\
Release Passport & release responsibility & What reviewed state and artifacts
were actually promoted and how that claim was verified \\
\bottomrule
\end{tabular}
\caption{Separate authorities composed without a universal project object.}
\label{tab:composition}
\end{table}

A Project Cut is a binding protocol, not a fourth fact engine. Its dependency
direction is one-way:

\begin{equation}
\begin{split}
\text{source}_n + \mathrm{Atlas}_n &\rightarrow \mathrm{Episodes}_{n+1}\\
\text{source}_{n+1} + \mathrm{Episodes}_{n+1}
  &\rightarrow \mathrm{Atlas}_{n+1}\\
\text{source projection} + \mathrm{Atlas}_{n+1}
  + \mathrm{EpisodeDelta}_{n+1} &\rightarrow \mathrm{ProjectCut}_{n+1}\\
\mathrm{ProjectCut}_{n+1} &\rightarrow \text{Git publication}
  \rightarrow \text{Release Passport}.
\end{split}
\end{equation}

Git ancestry does not become causal truth; Episode does not become complete
project knowledge; Atlas does not become runtime history; Passport does not
decide whether the promoted candidate is valuable. The binding preserves
separate roots and explicit responsibility. This is a concrete KFD-1/2 design:
facts remain non-drifting inside their authorities, and cross-authority trust
starts from inspectable bindings rather than name equality.

The current Project Cut decision is accepted but only partially implemented.
Its schema, provider equivalence, settlement, authority cutover, and full
qualification remain staged. It is used here as a composition design and
falsifiable implementation plan, not as completed evidence.

\subsection{The learning loop}

The composition supports a larger loop:

\begin{equation}
\mathrm{Atlas}_n
\rightarrow \text{guides work}
\rightarrow \mathrm{Episodes}_{n+1}
\rightarrow \text{reality pressure}
\rightarrow \text{assessment}
\rightarrow \mathrm{Atlas}_{n+1}.
\end{equation}

In KFD-5, grounded human judgment and agent reasoning interpret the pressure.
In the KFD-6 experiment, an agent mines many Episodes and proposes a candidate
closure, while held-out evidence and separate promotion authority prevent the
loop from certifying itself. Release Passport records what version of the loop
was actually promoted. The result is not a claim of automatic evolution; it is
an auditable path by which a system can learn without replacing external
consequence with generated narrative.

\subsection{Claim ladder}

Evaluation should separate increasingly strong claims:

\begin{description}[leftmargin=*,style=nextline]
  \item[C0: coherent declaration.]
    The candidate has explicit identity, boundary, authority, lifecycle,
    operations, fact sources, observer, alternatives, and falsifiers.
  \item[C1: reproducible artifact.]
    Independent readers can verify the candidate record, roots, and declared
    evidence from public or exported material.
  \item[C2: operational closure.]
    The object supports its minimum operations under interruption, drift, and
    recovery rather than only validating a schema.
  \item[C3: cross-participant reconstruction.]
    A new human or agent can continue work without private oral context and
    without silently creating a second authority.
  \item[C4: generalization.]
    The candidate reduces reconstruction or coordination cost across projects,
    participants, or providers while preserving local facts and falsifiers.
\end{description}

Historical importance, ecosystem lock-in, or universal applicability lies
beyond C4 and cannot be established by the originating team alone.

\subsection{Evaluation protocols}

\paragraph{Deletion and fuse.}
Remove the candidate from a controlled workflow. Measure which coordination
must be reconstructed, which facts become ambiguous, and whether a narrower
alternative closes the gap. Conversely, test whether the candidate genuinely
fuses previously separate mechanisms or merely renames them.

\paragraph{Cold-start recovery.}
Give a fresh participant only the declared artifacts and roots. Measure time,
private-context requests, unsupported assumptions, and correctness of the
reconstructed next action. Delete caches and projections to test stage-zero
recovery from authoritative material.

\paragraph{Drift and fault injection.}
Change a welded surface without its compatibility boundary; omit an Episode;
present a stale Atlas; alter an artifact after release; or mismatch a receipt
and cut. The system should fail visibly or report degraded/unverifiable rather
than synthesize green state.

\paragraph{Held-out and prospective evidence.}
For KFD-6, separate discovery Episodes from evaluation Episodes. Include cases
where the same internal pattern occurs without the proposed boundary. Run the
candidate in shadow mode, then bounded prospective interventions. Measure
compression, prediction, transfer, disconfirmation, and intervention outcome
against retrieval, summarization, and no-new-primitive baselines.

\paragraph{Independent assessment.}
The candidate generator cannot be the only verifier. A separate evaluator,
policy, evidence cut, or accountable participant must be able to reject the
claim. Promotion authority remains separate from discovery authority and
proportionate to consequence.

\subsection{Current evidence status}

At the July 15, 2026 evidence cut, KFD-1 through KFD-5 are active and KFD-6 is
draft. The KFD package publishes machine-readable decision and schema surfaces.
Xinfa Atlas identity and bounded projections are implemented with retained
qualification references. Buildchain Release Passport operates in stable
releases. Episode has a substantial product and runtime architecture but mixed
implementation status across its identity and portability contracts. Project
Cut is accepted and partial.

This supports C0--C2 unevenly across the cases. It does not yet support a
comparative C3 study, a C4 generalization claim, or activation of KFD-6. Those
gaps define the next research work rather than being hidden by the conceptual
coherence of the model.

\section{Related Work}
\label{sec:related}

\subsection{Human augmentation}

Engelbart framed computing as augmenting human intellect rather than merely
automating tasks \cite{engelbart1962augmenting}. This paper extends that
augmentation stance to real-world work with agents. The product should not
only automate actions; it should make the shared work environment more
legible for human and agent participants.

\subsection{Situated action and coordination}

Suchman's work on situated action argues that plans and situated practice have
a complex relationship \cite{suchman1987plans}. Winograd and Flores emphasize
language, commitments, and coordination in work \cite{winograd1986understanding}.
The KFD foundation shares the view that work is not reducible to static
instructions. It differs by focusing on the engineering of fact sources,
trust, and collaboration surfaces for agent-era products.

\subsection{Distributed systems and event evidence}

Lamport's work on time and ordering shows that physical time is not enough for
causal reasoning in distributed systems \cite{lamport1978time}. Event sourcing
treats current state as derived from recorded events \cite{fowler2005eventsourcing}.
Distributed tracing records runtime relationships across components
\cite{opentelemetry}. These mechanisms provide important evidence, but the KFD
foundation addresses a broader product trust question: how facts become
non-drifting, how claims become inspectable, and how participants cooperate
through trusted value.

\subsection{Software supply chain and provenance}

Software supply-chain mechanisms such as provenance attestations and
reproducible builds aim to make artifacts more trustworthy \cite{slsa}.
KFD-2 is compatible with this direction but generalizes beyond release
artifacts. The same trust assessment structure can apply to configs, schemas,
collaboration interfaces, timeline views, runtime facts, and agent control
surfaces.

\subsection{Agent tooling}

Current agent tooling often emphasizes tracing, evaluation, prompt
management, or workflow execution. These are valuable subsystems. The KFD
foundation argues for a higher layer: real-world agent work needs a shared
fact and trust model before tracing or control can become durable cooperation.

\section{Discussion and Limitations}
\label{sec:discussion}

\subsection{A systems foundation, not a moral ranking}

The words \emph{facts}, \emph{trust}, and \emph{cooperation} carry moral and
political meanings. This paper does not rank participants by possession of
those virtues. It assigns operational responsibilities to a system: preserve
the reference, expose the warrant, and make the basis of coordinated action
inspectable. A participant may be honest while using a drifting contract; a
reputable institution may make an under-evidenced claim; competitors may
interoperate reliably without mutual affection.

\subsection{Facts are admitted, not total}

KFD-1 does not turn a database, repository, or ledger into reality. Every fact
cut has an observer, source, authority boundary, and omissions. Stability can
preserve a mistaken model. The recursive kernel matters because contrary
consequence can be retained without rewriting the earlier cut, making
correction possible. KFD-4 makes the perspective boundary more explicit; it
does not make all perspectives equally supported.

\subsection{Trust is bounded, not accumulated prestige}

Trust assessments are purpose-specific and revisable. A system may be trusted
to preserve bytes but not to interpret their value; an agent may be trusted to
run a bounded transformation but not to approve its deployment. Aggregating
these assessments into a universal trust score would discard the purpose,
evidence, and responsibility boundaries on which KFD-2 depends.

\subsection{Cooperation is not harmony or permissiveness}

Participants may disagree, compete, refuse, or operate under strict safety and
legal constraints. KFD-3 applies where a system asks reasoning participants to
coordinate. It requires participant-relevant value and constraints to be
inspectable; it does not guarantee equal bargaining power or eliminate
coercive conditions outside the system boundary. A complete evaluation must
therefore examine meaningful alternatives and affected non-participants, not
only interface transparency.

\subsection{Minimal roles are not a complete ontology}

The removal arguments establish that the three roles are non-redundant within
the proposed model. They do not prove that the English terms are uniquely
minimal or that every domain must expose them as first-class objects. A system
that preserves equivalent continuity through another decomposition is a valid
counterexample to vocabulary-level necessity. Security, incentives, law,
embodiment, resource constraints, institutional history, and domain semantics
remain outside the triad's complete account.

\subsection{Originating-ecosystem bias}

The model was developed inside the Kungfu ecosystem and is reflected in its
product, release, and agent-work designs. That coherence is implementation
evidence but also a source of overfitting. The cross-session case is a design
trace, not proof that the same decomposition minimizes cost in healthcare,
public administration, scientific collaboration, or other domains.
Independent adopters may find a narrower kernel, different boundaries, or
social practices that outperform formalization.

\subsection{Formalization has a cost}

Fact cuts, evidence bindings, responsibility assignments, and cooperation
records move reconstruction cost from participants into infrastructure. The
move is justified only where the recurring cost, consequence, or duration of
the work exceeds the cost of maintaining the mechanism. A short-lived local
task may need no formal KFD surface. The foundation is load-bearing where the
work must cross a boundary and still remain the same work.

\section{Conclusion}

Real-world agent work changes the product problem. A product is no longer only
a tool used by a human or an API called by automation. It can become a shared
work environment where humans and agents observe facts, evaluate claims,
choose paths, accept constraints, and leave records for future recovery.

This paper proposed a compact foundation model for that environment:
non-drifting facts, inspectable trust, and trusted value. The three layers are
ordered and mutually supporting. Facts that drift cannot support trust. Trust
that is not inspectable cannot support responsible control. Value that is not
trustable cannot support voluntary cooperation.

The next step is evidence. The KFD package, Buildchain, and Kungfu provide a
path for turning the model into public-safe artifacts, release gates, local
fact surfaces, and real recovery cases. If those systems can bear the model
under real complexity, the foundation becomes more than a philosophy. It
becomes an engineering discipline for cooperating with agents in reality.


\bibliographystyle{plain}
\bibliography{references}

\end{document}
